\documentclass[10pt,a4paper]{article}

\usepackage[a4paper,left=18mm,right=18mm,top=18mm,bottom=18mm]{geometry}
\usepackage[T1]{fontenc}
\usepackage[utf8]{inputenc}
\usepackage{helvet}
\renewcommand{\familydefault}{\sfdefault}
\usepackage{microtype}
\usepackage{booktabs}
\usepackage{tabularx}
\usepackage{array}
\usepackage{graphicx}
\usepackage{xcolor}
\usepackage{enumitem}
\usepackage{parskip}
\usepackage{titlesec}
\usepackage{fancyhdr}
\usepackage{caption}
\usepackage{float}
\usepackage{hyperref}

\definecolor{ThesisGreen}{HTML}{184B3A}
\definecolor{ThesisBlue}{HTML}{547797}
\definecolor{ThesisBrick}{HTML}{B5675E}
\definecolor{ThesisGold}{HTML}{C78B34}
\definecolor{ThesisInk}{HTML}{26323B}
\definecolor{ThesisGrey}{HTML}{6B7280}
\definecolor{ThesisLight}{HTML}{EEF3F0}
\definecolor{ThesisRule}{HTML}{D7DEE3}

\hypersetup{
  colorlinks=true,
  linkcolor=ThesisGreen,
  urlcolor=ThesisBlue,
  citecolor=ThesisGreen,
  pdfauthor={Tim Kuhleis},
  pdftitle={Full-cohort acquisition DiD: design, results, and decisions}
}

\titleformat{\section}
  {\Large\bfseries\color{ThesisGreen}}
  {\thesection}{0.65em}{}
  [\vspace{-0.25em}\color{ThesisRule}\titlerule]
\titleformat{\subsection}
  {\large\bfseries\color{ThesisInk}}
  {\thesubsection}{0.55em}{}
\titlespacing*{\section}{0pt}{1.25em}{0.55em}
\titlespacing*{\subsection}{0pt}{1em}{0.35em}

\captionsetup{
  font=small,
  labelfont={bf,color=ThesisGreen},
  textfont={color=ThesisInk},
  skip=5pt
}
\setlist[itemize]{leftmargin=1.35em,itemsep=0.25em,topsep=0.35em}
\setlist[enumerate]{leftmargin=1.5em,itemsep=0.3em,topsep=0.35em}
\renewcommand{\arraystretch}{1.12}

\pagestyle{fancy}
\fancyhf{}
\lhead{\color{ThesisGrey}\small Full-cohort acquisition DiD}
\rhead{\color{ThesisGrey}\small Supervisor results package}
\cfoot{\color{ThesisGrey}\small \thepage}
\renewcommand{\headrulewidth}{0.3pt}
\renewcommand{\headrule}{\hbox to\headwidth{\color{ThesisRule}\leaders\hrule height \headrulewidth\hfill}}

\newcolumntype{Y}{>{\raggedright\arraybackslash}X}
\newcommand{\keybox}[1]{%
  \noindent\colorbox{ThesisLight}{%
    \parbox{\dimexpr\linewidth-2\fboxsep\relax}{\color{ThesisInk}#1}}}
\newcommand{\decisionbox}[1]{%
  \noindent\fcolorbox{ThesisGold}{white}{%
    \parbox{\dimexpr\linewidth-2\fboxsep-2\fboxrule\relax}{#1}}}
\newcommand{\sourceNote}[1]{%
  \par\smallskip{\footnotesize\color{ThesisGrey}\textit{Notes:} #1}}

\input{generated_values.tex}

\begin{document}

\begin{titlepage}
  \thispagestyle{empty}
  \vspace*{8mm}
  {\color{ThesisGreen}\rule{\textwidth}{2.2pt}}\par
  \vspace{10mm}
  {\fontsize{27}{31}\selectfont\bfseries\color{ThesisInk}
   Pharmaceutical acquisitions and inventor patenting\par}
  \vspace{4mm}
  {\Large\color{ThesisGreen}
   Full-cohort DiD: design, certified results, and decisions\par}
  \vspace{12mm}
  {\large Tim Kuhleis\par}
  \vspace{2mm}
  {\color{ThesisGrey}Supervisor results package \textbullet\ \BuildDate\par}
  \vfill
  {\color{ThesisGreen}\rule{\textwidth}{2.2pt}}\par
\end{titlepage}

\section{Proposed thesis structure}

\begin{table}[H]
\centering
\small
\begin{tabularx}{\textwidth}{@{}Y r Y@{}}
\toprule
Section & Pages & Main content \\
\midrule
Abstract & 0.5 & Question, design, result, and qualification \\
1. Introduction & 2.0 & Motivation, contribution, result, and roadmap \\
2. Background, Literature, and Hypotheses & 4.0 & Pharma acquisitions,
patent measurement, mechanisms, and hypotheses \\
3. Data and Sample Construction & 4.5 & Data, deal and inventor construction,
outcomes, and descriptives \\
4. Empirical Design & 5.0 & Placebo cohorts, local support, entropy
balancing, LOYO, and inference \\
5. Main Results: Full Target-Inventor Cohort & 5.5 & Balance, quantity,
LOYO, margins, and secondary outcomes \\
6. Stayers, Selection, and Individual Heterogeneity & 5.5 & Initially
retained definition, separate balance, selection, and bounds \\
7. Technological Relatedness and Mechanisms & 4.0 & DealSim, TechDrift,
relatedness heterogeneity, and mechanisms \\
8. Robustness and Interpretation & 2.5 & Inference, support, timing,
censoring, and scope \\
9. Conclusion & 1.5 & Findings, interpretation, policy relevance, and limits \\
\midrule
\textbf{Total} & \textbf{35.0} & \\
\bottomrule
\end{tabularx}
\end{table}

\decisionbox{\textbf{Discussion question 1.} Does this proposed 35-page
structure and page allocation seem appropriate, including the balance between
the full-cohort, initially retained-inventor, and mechanism results?}

\decisionbox{\textbf{Discussion question 2.} Keeping the full cohort,
initially retained inventors, and a full technological-relatedness section
risks producing three main-results papers inside one thesis. Should Section 7
remain four pages, be compressed to roughly 2--2.5 pages, or move largely to
the appendix?}

\clearpage
\section*{Executive summary}

\begin{enumerate}
  \item \textbf{The conventional estimator is not credible as-is.}
  A standard Callaway--Sant'Anna (2021) DiD with all not-yet-treated
  inventors, no placebo treatment clock, and no matching exhibits strong
  pre-treatment differences. In this sample, the standard CS implementation
  is not a usable causal design without further counterfactual construction.

  \item \textbf{Recruitment timing explains the mechanical hump, but not all
  lifecycle differences.} Inventors qualify through a recent patent. The raw
  peak moves with each inventor's qualifying year. Assigning controls placebo
  deal years aligns this clock, yet the unmatched placebo DiD still rejects
  parallel pre-trends (\(p<0.001\)).

  \item \textbf{The preferred design must create comparable pre-treatment
  populations.} Donors first satisfy prospectively frozen acquisition-clean
  firm restrictions, IPC4 technological support, firm and inventor calipers,
  and control-diversification rules. I then use entropy balancing, following
  \href{https://doi.org/10.1093/pan/mpr025}{Hainmueller (2012)}, to reweight
  admissible controls so their five-year pre-deal patent paths, inventor
  characteristics, and firm characteristics resemble those of treated
  inventors within acquisition cohorts.

  \item \textbf{The matched estimate is economically large and statistically
  precise.} The headline ATT is \(\HeadlineATT\) patents per inventor-year
  (95\% CI \([\HeadlineLow,\HeadlineHigh]\), \(p=\HeadlineP\)). This equals
  \(\CumulativeMagnitude\) fewer patents over five years, a
  \(\PreDeclinePct\%\) decline relative to average pre-acquisition output and
  \(\CounterfactualDeclinePct\%\) of estimated post-acquisition
  counterfactual patenting.

  \item \textbf{The result survives every leave-one-pre-year-out design.}
  All five LOYO estimates are negative and statistically significant. Among
  directly comparable specifications, the largest point-estimate deviation
  is only \(\MaxLOYOPct\%\) of the headline magnitude. Positive held-out gaps
  at \(t=-3\) and \(t=-2\) nevertheless make us wary about parallel trends
  and limit the strength of the causal interpretation.

  \item \textbf{The robust conclusion is about quantity.} Approximately
  \(\ExtensiveSharePct\%\) of the accounting decline operates through a lower
  probability of patenting and \(\IntensiveSharePct\%\) through fewer patents
  in active years. Patent quality, measured by the OECD PQII composite index,
  falls in the full sample; TechDrift is uninformative; and forward citations
  rise only in the full cohort window. These composition outcomes remain
  secondary.
\end{enumerate}

\keybox{\textbf{Bottom line.} The full-cohort matched ATT is the closest
supported approximation to the causal effect available in these data. Its
magnitude is strikingly stable across recruitment, censoring, and LOYO
diagnostics. The remaining held-out pre-period humps should be reported
directly as an identification concern, not hidden by the exact balance
of the headline specification.}

\section{Data and sample}

The analysis starts from 513 transactions in the Cassi--Ornaghi acquisition
and patent data. The design needs five patent years before and five years
after each transaction. Given the 1988--2015 patent window, the estimating
sample is therefore restricted to 1994--2010 and contains
\(\EligibleDeals\) acquisitions and \(\EligibleInventors\) target inventors
who patented with a target during the five years before its deal.

Patent applications decline sharply in the final source years, which raises a
separate concern that right-censoring may depress measured post-deal output.
The buffered 1994--2008 cohort reduces exposure to those years and provides a
direct censoring check. Descriptive tables are imported from the certified
stage-14 data-section package.

\input{assets/table1_sample_construction.tex}

\input{assets/table3_overall_descriptives.tex}

The target inventor is relatively early in the observed patenting career:
median career age is three years, median pre-deal patent stock is one, and
median focal-group exclusivity equals one. This makes lifecycle and
recruitment-clock comparability central rather than incidental. Treatment is
also concentrated: the five largest deals account for 29.6\% of target
inventors and the ten largest for 42.7\%. The inventor-share HHI corresponds
to 38.2 equally sized deals. This motivates deal-level inference and explicit
reporting of effective treatment variation.

\section{Why the standard DiD fails, and what replaces it}

The first panel of Figure~\ref{fig:designsequence} applies the conventional
Callaway--Sant'Anna (2021) approach: all not-yet-treated inventors serve as
controls, with no placebo timing and no matching. The pre-period divergence is
large. Panel B then shows why the raw hump is partly mechanical: each
qualifying-year group peaks in its own qualifying year. Panel C aligns treated
and control recruitment clocks by assigning placebo deal years. This removes
the timing mismatch but leaves different lifecycle paths. Consequently, the
unmatched placebo ATT is close to zero
(\(\RawPlaceboATT\), \(p=\RawPlaceboP\)), while its joint pre-trend test still
rejects (\(p<0.001\); Panel D).

\begin{figure}[H]
  \centering
  \includegraphics[width=\textwidth]{assets/figure_d1_design_sequence.pdf}
  \caption{From the conventional CS design to the matched placebo design}
  \label{fig:designsequence}
  \sourceNote{Panel A uses all not-yet-treated inventors without placebo
  cohorts or matching. Panels B--D diagnose the recruitment clock and
  residual lifecycle differences. Confidence intervals are 95\%.}
\end{figure}

The preferred design therefore uses placebo cohorts, local support, and
entropy balancing. Following \href{https://doi.org/10.1093/pan/mpr025}
{Hainmueller (2012)}, entropy balancing is not a cosmetic adjustment: it
reweights admissible controls to reproduce the treated cohort's specified
pre-treatment moments and thereby adjusts the comparison for observed
lifecycle differences.

\subsection{Prospectively frozen support restrictions}

Support is constructed separately within each acquisition cohort \(g\):

\begin{enumerate}
  \item \textbf{U2-clean donor firms.} A control firm has no target event on or
  before \(g+5\) and no acquirer event during \(g-5,\ldots,g+5\). Control
  inventors exposed to a target event on or before \(g+5\) are excluded.
  \item \textbf{Firm support.} Candidate firms share IPC4 technology with the
  target and satisfy the frozen Stage-1 caliper of 2.0. The nearest 50
  admissible firms are ranked using IPC4 cosine similarity, five-year firm
  patent stock, and five-year inventor count. Firm patent trajectory remains
  an exact balance moment.
  \item \textbf{Inventor support.} Inventor pairs share IPC4 support and satisfy
  the Stage-2 caliper of 1.5 on five-year output, the annual patent trajectory,
  career age, and IPC4 cosine distance. Each supported treated inventor has
  at least three admissible controls from at least two firms.
\end{enumerate}

Within this support, the final design balances annual patent counts and
active-patenting indicators for \(t=-5,\ldots,-1\), career age,
focal-group exclusivity, firm patent stock, firm inventor count, and firm
patent trajectory.

\begin{table}[H]
\centering
\caption{Support, balance, and effective information}
\label{tab:designsummary}
\input{generated_design_table.tex}
\end{table}

\begin{figure}[H]
  \centering
  \includegraphics[width=0.80\textwidth]{assets/figure_d2_love_plot.pdf}
  \caption{Balance before and after entropy balancing}
  \label{fig:loveplot}
  \sourceNote{Each point is the largest absolute standardized mean difference
  across the 17 acquisition cohorts. Post-deal outcomes are not balance
  targets.}
\end{figure}

\section{Main quantity result and identification diagnostics}

The preferred design estimates \(\HeadlineATT\) patents per inventor-year over
\(t=+1,\ldots,+5\). Accumulated over five years, this is
\(\CumulativeMagnitude\) fewer patents per inventor. The effect is
\(\PreDeclinePct\%\) of average pre-treatment patenting and
\(\CounterfactualDeclinePct\%\) of the post-treatment output predicted for
treated inventors under the matched counterfactual.

Patent counts fall sharply at the end of the source data: from 54,850
applications in 2013 to 28,096 in 2014 and 1,143 in 2015. The 1994--2008
companion therefore checks whether late-cohort right-censoring creates the
effect. It produces \(\BufferedATT\) patents per inventor-year
(\(p=\BufferedP\)), nearly identical to the headline. Censoring is present in
the source data but does not explain the quantity result.

\subsection{Leave-one-pre-year-out placebo designs}

Balancing every displayed pre-treatment outcome makes the headline leads zero
by construction. Hainmueller's entropy-balancing explainer recommends leaving
at least one pre-period outside the constraints when a genuine placebo is
desired. The complete LOYO grid does this one year at a time.

\begin{figure}[H]
  \centering
  \includegraphics[width=0.94\textwidth]{assets/figure_d3_patent_count_loyo.pdf}
  \caption{The patent-count result survives every held-out pre-year}
  \label{fig:loyo}
  \sourceNote{Average annual ATT over \(t=1,\ldots,5\). Bars show 95\%
  deal-level wild-bootstrap confidence intervals. The \(t=-1\) holdout uses
  \(t=-4\) as its reference and is therefore not point-comparable to the
  other rows.}
\end{figure}

\input{generated_loyo_table.tex}

Every LOYO effect is negative and statistically significant. Among the
directly comparable \(t=-5,\ldots,-2\) holdouts, the estimate ranges from
\(\LOYOMin\) to \(\LOYOMax\); the largest deviation is only
\(\MaxLOYOPct\%\) of the headline magnitude. The result therefore survives
the specifications in which the pre-treatment humps are allowed to appear.

The held-out gaps at \(t=-3\) and \(t=-2\) are positive and statistically
different from zero. This means the full-sample exact balance cannot, by
itself, prove parallel trends. A transparent additive calibration implies a
sign-breakdown multiple of \(\BreakdownMultiple\): an unobserved post-period
violation would have to equal \(\BreakdownMultiple\) times the largest
held-out discrepancy to offset the point estimate. This is a descriptive
sensitivity statistic, not a formal HonestDiD confidence set.

The early-recruitment design inspired by
\href{https://www.nature.com/articles/s41599-025-04894-w}
{Verginer and Riccaboni} separates recruitment at \(t=-7,-6\) from the later
pre-treatment observation period. Where feasible, its estimates remain
negative. It is nevertheless more selective, changes the estimand toward
experienced long-tenure inventors, and makes strict firm-and-inventor balance
infeasible in three cohorts. It is reported in the appendix as
established-inventor heterogeneity, not as a superior replacement design.

\section{What drives the patent decline?}

The matched post-acquisition counterfactual is \(\ControlPostMean\) patents
per inventor-year, compared with \(\TreatedPostMean\) among treated inventors.
Their difference can be separated into the probability of patenting in a year
and output in active inventor-years.

\input{generated_margin_table.tex}

The probability of patenting falls from \(\ControlActivePct\%\) to
\(\TreatedActivePct\%\), a \(\ActiveDropPP\)-percentage-point difference.
This extensive margin accounts for \(\ExtensiveSharePct\%\) of the total.
The remaining \(\IntensiveSharePct\%\) is associated with fewer patents in
active inventor-years: \(\ControlIntensive\) under the counterfactual versus
\(\TreatedIntensive\) for treated inventors.

The latter comparison is not a causal ATT conditional on activity, because
acquisition can itself determine who remains active. Its substantive
interpretation is deferred to the separately balanced analysis of inventors
initially retained in the merged entity.

\subsection{Predetermined inventor heterogeneity}

\input{generated_full_heterogeneity_table.tex}

The quantity decline is larger among inventors who were more productive
before acquisition, repeatedly patented with a persistent collaborator, or
had greater pre-deal technological overlap with the acquirer. The productivity
gradient is not only an absolute-stock arithmetic result: the predicted loss
increases from approximately 4.4\% of pre-deal output at the treated
productivity p25 to 14.7\% at p75. The persistent-team result is consistent
with disruption of relationship-specific knowledge, while the negative
TechFit gradient is consistent with redundant programs being rationalized.
These are moderator patterns, not direct tests that uniquely identify either
mechanism. Career age does not materially moderate the patent-count effect,
and none of the active-patenting contrasts is conventionally significant.

\begin{figure}[H]
  \centering
  \includegraphics[width=0.88\textwidth]
    {assets/figure_vr_heterogeneity_focal.png}
  \caption{Full-cohort heterogeneity across predetermined inventor traits}
  \label{fig:full-heterogeneity}
  \sourceNote{Points compare the frozen high-versus-low moderator values.
  All eight predeclared contrasts are shown.}
\end{figure}

\section{Secondary outcomes}

\input{generated_secondary_table.tex}

\begin{figure}[H]
  \centering
  \includegraphics[width=0.92\textwidth]{assets/figure_d5a_all_outcomes_event_study_full.pdf}
  \caption{Matched event studies for all outcomes, 1994--2010 cohorts}
  \label{fig:secondaryfull}
  \sourceNote{Patent count and active patenting are balance targets, so their
  headline pre-period coefficients are zero by construction; LOYO supplies
  their genuine placebo checks. PQII, TechDrift, and forward citations were
  not balance targets and retain genuine pre-period diagnostics.}
\end{figure}

The secondary results are deliberately not forced into a single mechanism:

\begin{itemize}
  \item \textbf{Probability of patenting} declines by about 1.8 percentage
  points. The full-sample average narrowly misses 5\% significance
  (\(p=\ActiveP\)), although several individual post years are negative and
  significant under two-way clustering. Conservative deal-level inference
  and only \(\EffectiveDeals\) effective treated deals limit precision.
  \item \textbf{OECD PQII} is negative in the full sample
  (\(p=\PQIIP\)) and has reassuring unbalanced pre-period evidence. It is
  smaller and imprecise in the buffered sample, so incomplete linkage and
  late-sample coverage make it secondary.
  \item \textbf{Five-year forward citations} are positive in the full sample
  (\(p=\CitationP\)) but smaller and insignificant in the buffered sample.
  This is compatible with selective filing among the patents that remain; it
  does not reverse the quantity result.
  \item \textbf{TechDrift} is small, imprecise, and fails its unbalanced joint
  pre-trend diagnostic. The data do not identify a change in technological
  distance. This is an informative null and should be kept as a secondary
  outcome rather than optimized for significance.
\end{itemize}

\section{Interpretation, limits, and thesis framing}

The evidence supports a precise but qualified conclusion. A raw post-deal
decline is not credible evidence because recruitment timing and inventor
lifecycles generate similar shapes in untreated inventors. Once those clocks
are aligned and observable five-year histories are balanced within clean
local support, treated inventors produce materially fewer patents. The
quantity effect survives every LOYO specification, the buffered cohort
restriction, and the feasible early-recruitment design.

The positive \(t=-3\) and \(t=-2\) LOYO gaps make us wary about parallel
trends and limit the strength of the causal interpretation. They prevent an
assumption-free or unqualified causal claim. At the same time, their presence
does not explain away the post-treatment magnitude. The appropriate
interpretation is that the matched ATT is the closest supported causal
approximation available in these data, with sensitivity stated transparently.

\keybox{\textbf{Thesis framing.} Raw inventor patenting declines after
acquisition, but untreated inventors exhibit similar lifecycle dynamics and
differ from treated inventors before the deal. I therefore restrict donors to
clean local firm and IPC4 support and use entropy balancing to align complete
five-year pre-acquisition patent paths and lifecycle characteristics. The
matched design estimates \(\AnnualMagnitude\) fewer patents per inventor-year,
equivalent to a \(\PreDeclinePct\%\) decline relative to average
pre-acquisition output and \(\CounterfactualDeclinePct\%\) of estimated
post-acquisition counterfactual patenting. The estimate remains negative and
statistically significant in every leave-one-pre-year-out specification, with
a maximum point-estimate deviation of \(\MaxLOYOPct\%\). The held-out
pre-period humps prevent an unqualified causal interpretation, but neither the
lifecycle nor alternative-recruitment diagnostics explains away the post
effect.}

\decisionbox{\textbf{Discussion question 3.} Is the combination of clean local
support, exact entropy balance, the LOYO diagnostics, and explicit
parallel-trends qualification sufficiently credible for the full-cohort
quantity result to anchor the thesis?}

\clearpage
\section{Initially retained inventors}

I classify an inventor as \emph{initially retained} when their first observed
patent in \(t=+1,\ldots,+5\) is assigned to the target, acquirer, or resulting
corporate group. The acquisition year is excluded because within-year timing
is ambiguous. Controls satisfy the symmetric rule: their first post-event
patent must be assigned to their clean control firm. This definition measures
initial patent-based retention, not continuous employment.

The stayer design is rebuilt separately rather than applying the full-cohort
weights after conditioning on a post-treatment outcome. It uses the same clean
firm and IPC4 support logic and entropy-balances the annual patent-count and
active-patenting histories over \(t=-5,\ldots,-1\), alongside inventor and firm
characteristics.

\input{generated_stayer_table.tex}

The point estimate is \(\StayerATT\) patents per inventor-year, or
\(\StayerCumulative\) over five post-acquisition years. This equals
\(\StayerCounterfactualPct\%\) of matched counterfactual post-deal output and
\(\StayerPrePct\%\) of average pre-deal output. The buffered estimate is
similar (\(\StayerBufferedATT\); deal-wild sensitivity
\(p=\StayerBufferedP\)), so late-sample censoring does not explain the result.
The frozen primary retained result uses the two-way deal/inventor interval,
which excludes zero. Deal-level wild inference is reported transparently as a
cluster-sensitivity companion and includes zero.

This result is economically large but only suggestive. First, initial
retention is determined after acquisition, so the estimate is a
selected-group contrast and not an effect for an always-retained causal
stratum. Second, the genuinely held-out \(t=-3\) and \(t=-2\) gaps are
\(\StayerMThreeGap\) and \(\StayerMTwoGap\), respectively. They make the
parallel-trends assumption less reassuring and prevent an unqualified causal
interpretation. The consistent negative result across inference and
censoring variants is nevertheless useful evidence that the full-cohort
decline is not confined to inventors who immediately cease patenting for the
combined group.

\subsection{Selection, margins, and heterogeneity}

\input{generated_stayer_selection_table.tex}

Initially retained inventors were already positively selected on pre-deal
patenting: their average five-year patent stock and number of active pre-deal
years exceed those of both leavers and inventors with no observed post-deal
patent. This does not invalidate the conditional result, but it confirms that
retention is post-treatment selected and that the estimate cannot be
generalized to all treated inventors or an unobserved always-retained stratum.

\input{generated_stayer_decomposition_table.tex}

The P8 accounting decomposition assigns 27.3\% of the initially retained
decline to an earlier end of observed patenting, 5.4\% to fewer active years
among patenting survivors, and 67.3\% to fewer patents per active year. The
cessation component is partly mechanical because initial retention is defined
through patenting in \(t=+1,\ldots,+5\). The full-cohort cessation result
therefore carries the stronger substantive claim.

The predetermined retained-sample heterogeneity estimates are available in
the accompanying appendix artifact but are too imprecise to carry a separate
main-text conclusion.

\begin{figure}[H]
\centering
\includegraphics[width=0.66\textwidth]
  {assets/figure_s4_patent_paths_and_event_study.pdf}
\caption{Initially retained inventors: weighted patent paths and dynamic
contrasts}
\label{fig:stayer-paths}
\sourceNote{Panel A shows weighted patent-count levels. Panel B reports
differences relative to \(t=-1\) with two-way deal/inventor intervals. The
headline ATT excludes the acquisition year and averages \(t=+1,\ldots,+5\).}
\end{figure}

\decisionbox{\textbf{Discussion question 4.} Is it appropriate to present the
initially retained estimate as economically important but suggestive
selected-group evidence, with the full-cohort ATT remaining the main result?}

\decisionbox{\textbf{Discussion question 5.} Is the positive \(t=0\) contrast
concerning, or is excluding the acquisition year appropriate because annual
patent data cannot order patents relative to deal completion and both groups
decline at \(t=0\), with the matched controls declining more?}

\clearpage
\appendix
\section{Compact appendix}

\subsection{Held-out pre-period gaps}
\input{generated_placebo_table.tex}

\subsection{Deal concentration}
\input{assets/table4_deal_inventor_concentration.tex}

\vfill
\begin{center}
\footnotesize\color{ThesisGrey}
Generated reproducibly from certified P5c/P6 outputs and stage-14 descriptive
tables. Build ID: \BuildID.
\end{center}

\end{document}
